\clearpage
\section{引言}

\subsection{项目背景与动机}
同一份数据由多个节点共同维护时，写入成功涉及两个问题：副本是否以相同顺序执行操作，以及客户端能否在节点切换后读到已经确认的结果。以键值写入为例，Leader 在本地追加命令后还需要把日志复制到足够多的节点。若此时 Leader 退出，剩余节点要从已有日志中选出能够继续服务的节点，并在原有提交位置上推进。

事务处理关注的是另一条顺序。一次操作可能同时修改多个键，读请求需要识别哪些版本对当前事务可见，进程恢复时还要依据持久记录识别提交结果。项目选择复制与事务恢复作为两条实现主线，并把操作历史和运行状态纳入系统，使一次写入的确认、应用与恢复可以沿记录追溯。

\subsection{系统目标}
系统面向键值访问与多键事务两类业务。键值集群承担副本管理，提供 Leader 选举、读写、日志追赶、快照与成员变更。事务服务提供版本读写、隔离级别选择、提交与中止、历史读取和崩溃恢复。两类服务分别维护复制索引和事务提交顺序，通过公共操作记录与观测接口接入验证模块和控制台。

\subsection{主要机制}
\textbf{复制状态机。} Raft 将选举、日志复制与提交组织在任期和多数派之上\cite{raft}。节点按已提交日志的顺序执行确定性命令，在相同的应用边界得到相同键值状态。线性一致性要求已完成操作的先后关系与实际时间相容\cite{linearizability}，系统在默认读路径中以当次多数派确认建立读屏障。

\textbf{多版本与持久化。} MVCC 将同一键的提交结果保存为版本链。读已提交在每次读取时选择当前已提交版本，快照隔离按事务开始时的提交位置选择版本\cite{isolation}。预写日志（WAL）先记录待物化动作和提交决定，恢复过程按日志顺序重建状态，并用补偿记录保存撤销进度\cite{aries}。

\textbf{运行记录与验证。} 客户端历史记录请求的调用、返回和可见结果。验证器由此检查单键操作顺序，由事务的读版本和写版本重建依赖，并在提交边界检查业务条件。观测控制台把这些结果同节点状态和运行事件一并呈现。
